\chapter [Dedica e aspirazione]{Versi di dedica e aspirazione}

\firstline{Iminā puññakammena upajjhāyā guṇuttarā}

\delegateSetUseNext

\begin{leader}
  [Ha꜓nda mayaṁ uddissanādhiṭṭhāna-gāthā꜓yo b꜕haṇāmase]
\end{leader}

\begin{english}
 [Cantiamo ora i versi di dedica e aspirazione]
\end{english}

% ? ref by Thai King
%\suttaref{Trad.}%
[Iminā puñña꜕kammena] u꜕pajjhāyā gu꜕ṇutta꜕rā\\

\begin{english}
[Attraverso il bene] che nasce dalla mia pratica,
\end{english}

Ācariyūpa꜕kārā ca꜕ mātāpitā ca꜕ ñāta꜕kā\\

  \begin{english}
Che i miei maestri spirituali di insuperabile virtù,
  \end{english}

Suriyo candimā rājā gu꜕ṇavantā na꜕rāpi꜕ ca꜕\\

  \begin{english}
 Che tutti i miei insegnanti, mia madre, mio padre e i miei parenti\\ 
 il Sole e la Luna\\
persone virtuose e guide del mondo\\
  \end{english} 

Brahma-mārā ca꜕ indā ca꜕ loka꜕pālā ca꜕ deva꜕tā\\

  \begin{english}
Possano gli dei supremi e le forze del male, gli spiriti guardiani della Terra e gli altri esseri celesti,
  \end{english}

Yamo mittā ma꜕nussā ca majjhattā veri꜕kāpi꜕ ca꜕\\
  \begin{english}
    Il Signore della Morte, persone amichevoli, indifferenti e ostili,
  \end{english}

Sa꜕bbe sattā sukhī hontu puññāni pa꜕ka꜕tāni꜕ me\\
  \begin{english}   
    Che tutti gli esseri ricevano i meriti della mia pratica.
  \end{english}

  \clearpage

Sukhañca tividhaṁ dentu꜕ khippaṁ pāpetha꜕ voma꜕taṁ\\

  \begin{english} 
Che presto giungano alla Triplice Felicità e realizzino il Senza-morte.
  \end{english}

Iminā puññakammena iminā uddi꜕ssena꜕ ca꜕\\
  \begin{english} 
Attraverso il bene che nasce dalla mia pratica e da questa dedica,
  \end{english}

Khipp'āhaṁ su꜕la꜕bhe ceva taṇhūpādāna꜕-cheda꜕naṁ\\

  \begin{english}
    Che possa io rapidamente e con facilità, estirpare ogni brama e attaccamento, 
  \end{english}

Ye santāne hīnā dhammā yāva꜕ nibbāna꜕to ma꜕maṁ\\

  \begin{english}   
Tutti gli stati mentali nocivi fino alla realizzazione del Nibbāna,
  \end{english}

Nassantu sabba꜕dā yeva yattha꜕ jāto bha꜕ve bha꜕ve\\

  \begin{english} 
    Possano dissolversi per sempre ovunque io rinasca in ogni esistenza
  \end{english}

Ujucittaṁ sa꜕ti꜕paññā sallekho vi꜕ri꜕yamhinā\\

  \begin{english} 
Possa mantenere una mente retta con consapevolezza e saggezza, austerità e vigore
  \end{english}

Mārā labhantu nokāsaṁ kātuñca vi꜕ri꜕yes꜕u me\\

  \begin{english} 
che le forze dell'illusione non trovino mai varco per ostacolare i miei sforzi
  \end{english}

Buddhādhipa꜕va꜕ro nātho dhammo nātho va꜕rutta꜕mo\\

  \begin{english}
Il Buddha è il mio sommo rifugio, Insuperabile è la protezione del Dhamma,
  \end{english}

Nātho pacceka꜕buddho ca꜕ saṅgho nāthotta꜕ro ma꜕maṁ\\

  \begin{english}
Il Buddha solitario è il mio nobile esempio,
Il Saṅgha è il mio più grande sostegno.
  \end{english}

Tesottamānubhāvena mārokāsaṁ la꜕bhantu꜕ mā

  \begin{english}
Armato di questo supremo potere che Mara non trovi mai varco.
  \end{english}

\clearpage

\chapter[Condivisione dei meriti]{Versi di condivisione dei meriti}

\firstline{Puññass'idāni katassa yān'aññāni katāni me}

\begin{leader}
  [Ha꜓nda mayaṁ sa꜕bba-patti-dāna-gāthā꜓yo bha꜕ṇāmase]
\end{leader}

   \begin{english}
  [Cantiamo ora i versi di condivisione dei meriti] 
\end{english}     


%\suttaref{Trad.}%
Puññass'i꜕dāni꜓ ka꜕ta꜕ssa yān'aññāni꜓ ka꜕tāni꜓ me\\
Tesa꜓ñca꜕ bhāgi꜕no ho꜓ntu꜕ sa꜕ttānantā꜕ppa꜕māṇa꜓kā

\begin{english}
 [Il merito] che ora ho compiuto\\
 e quello che ho accumulato in passato \\
Possa essere condiviso con tutte le creature,\\
 senza limiti in ogni direzione,
\end{english}

Ye pi꜕yā gu꜓ṇavantā ca꜕ mayhaṁ mātā-pi꜕tāda꜓yo\\
Diṭṭhā me cāpy꜕adiṭṭhā vā aññe majjh꜓att꜕a-veri꜓no

\begin{english}
 Persone care piene di bontà,\\
  a partire da mia madre e mio padre,\\
Esseri visti e non visti da me, indifferenti e ostili,
\end{english}

Sa꜕ttā tiṭṭha꜓nti꜕ lokasmiṁ te-bhummā ca꜕tu꜕-yoni꜓kā\\
Pañc'eka꜕-ca꜕tu꜕-vokārā sa꜓ṁsa꜕rantā bh꜕avābha꜕ve

\begin{english}
Esseri che dimorano nel mondo,\\
 nei tre piani di esistenza, dai quattro tipi di nascita,\\
Dotati di uno, quattro o cinque aggregati,\\
vagando da un'esistenza all'altra,
\end{english}

Ñātaṁ ye pa꜕tti꜓-dānam-me a꜕nu꜓modantu꜕ te sa꜕yaṁ\\
Ye c'imaṁ nappa꜕jānanti devā tesa꜓ṁ ni꜕veda꜕yuṁ

\clearpage

\begin{english}
 Coloro che sono a conoscenza della mia dedica,\\
  possano gioire di essa,\\
E per chi ancora non ne è a conoscenza\\
che gli déi gliela rendano nota.
\end{english}

Ma꜓yā dinnāna-puññānaṁ a꜕nu꜓moda꜕na-he꜓tu꜕nā\\
Sa꜕bbe sa꜕ttā sa꜕dā ho꜓ntu꜕ a꜕verā su꜕kh꜕a-jīvi꜓no\\
Kh꜓ema꜓ppa꜕dañca꜕ pa꜕ppontu꜕ tesā꜓sā꜓ si꜕jjha꜕taṁ su꜕bhā

\begin{english}
 E a causa del loro gioire per il merito donato da me,\\
Possano tutti gli esseri vivere sempre felici,\\
 liberi da ogni discordia\\
E possano loro realizzare il Nibbana\\
e le loro nobili aspirazioni.
\end{english}

\chapter[Mettā sutta]{Le parole del Buddha sulla gentilezza amorevole}

\firstline{Karaṇīyam-attha-kusalena}

\delegateSetUseNext

\begin{leader}
  [Ha꜓nda mayaṁ metta-sutta-gāthā꜓yo bha꜕ṇāmase]
\end{leader}

\begin{english}   
  [Cantiamo ora le parole del Buddha sulla gentilezza amorevole]
\end{english}



%\suttaref{Sn 1.8}%
[Karaṇīyam-attha-kusalena]\\
Yan-taṁ santaṁ padaṁ abhisamecca\\
Sakko ujū ca suhujū ca\\
Suvaco c'assa mudu anatimānī

\begin{english}  
[Questo dovrebbe fare] chi pratica il bene\\ 
e conosce il sentiero della pace:\\
Essere coraggioso, retto e onesto,\\
Semplice da ammonire, gentile e non arrogante
\end{english}

Santussako ca subharo ca\\
Appakicco ca sallahuka-vutti\\
Sant'indriyo ca nipako ca\\
Appagabbho kulesu ananugiddho

\begin{english}
  Contento, facile da aiutare,\\ 
  Non oppresso da impegni e di modi frugali,\\
  Calmo e prudente,\\
  Non altero o esigente,  
\end{english}

  \clearpage

Na ca khuddaṁ samācare kiñci\\
Yena viññū pare upavadeyyuṁ\\
Sukhino vā khemino hontu\\
Sabbe sattā bhavantu sukhit'attā

\begin{english}
 Incapace di fare ciò che il saggio poi disapprova.\\
 Che tutti gli esseri vivano felici e sicuri:
\end{english}

Ye keci pāṇa-bhūt'atthi\\
Tasā vā thāvarā vā anavasesā\\
Dīghā vā ye mahantā vā\\
Majjhimā rassakā aṇuka-thūlā

\begin{english}
Tutti, chiunque essi siano,\\
Risvegliati e non risvegliati,\\
Alti, medi e bassi, grandi e piccini
\end{english}


Diṭṭhā vā ye ca adiṭṭhā\\
Ye ca dūre vasanti avidūre\\
Bhūtā vā sambhavesī vā\\
Sabbe sattā bhavantu sukhit'attā

\begin{english}
  Visibili e non visibili,\\
  Vicini e lontani,\\
  Nati e non ancora nati.\\
  Che tutti gli esseri vivano felici!
\end{english}

\clearpage

Na paro paraṁ nikubbetha\\
Nātimaññetha katthaci naṁ kiñci\\
Byārosanā paṭighasaññā\\
Nāññam-aññassa dukkham-iccheyya

\begin{english}
  Che nessuno inganni l'altro,\\
  né lo disprezzi, né per odio o irritazione\\
  Desideri il suo male
\end{english}

Mātā yathā niyaṁ puttaṁ\\
Āyusā eka-puttam-anurakkhe\\
Evam'pi sabba-bhūtesu\\
Mānasam-bhāvaye aparimāṇaṁ

\begin{english}
  Come una madre proteggerebbe con la sua vita,\\
   suo figlio, il suo unico figlio\\
   Allo stesso modo si coltivi un cuore illimitato,\\
    verso tutti gli esseri viventi,
\end{english}

Mettañca sabba-lokasmiṁ\\
Mānasam-bhāvaye aparimāṇaṁ\\
Uddhaṁ adho ca tiriyañca\\
Asambādhaṁ averaṁ asapattaṁ

\begin{english}
Irradiando amore sull'universo intero;\\
In alto aldilà del cielo, in basso verso gli abissi,\\
Illimitatamente, in ogni direzione,\\
 liberi da odio e rancore.
\end{english}

\clearpage


Tiṭṭhañ-caraṁ nisinno vā\\
Sayāno vā yāvat'assa vigata-middho\\
Etaṁ satiṁ adhiṭṭheyya\\
Brahmam-etaṁ vihāraṁ idham-āhu

\begin{english}
Seduti, in piedi o camminando,\\
Distesi, esenti da torpore,\\
Sostenendo la pratica di Mettā,\\ 
Questa è la sublime dimora.
\end{english}

Diṭṭhiñca anupagamma\\
Sīlavā dassanena sampanno\\
Kāmesu vineyya gedhaṁ\\
Na hi jātu gabbha-seyyaṁ punaretī'ti

\begin{english}
  Il puro di cuore, dotato di retta visione,\\
Non legato ad opinioni,\\
Liberato da brame sensuali,\\
Non tornerà a nascere in questo mondo.
\end{english}

\clearpage


\chapter[Benessere universale]{Riflessione sul benessere universale}

\firstline{Ahaṁ sukhito homi}

\delegateSetUseNext

\begin{leader}
[Ha꜓nda mayam mettāpharaṇaṁ ka꜕romase]
\end{leader}

\begin{english}
[Cantiamo ora la riflessione sul benessere universale.]
\end{english}


%\suttaref{Trad.}%
[Aha꜓ṁ sukhito ho꜓mi]\\
  \begin{english}
[Che io possa essere felice]
\end{english}

Niddukkho ho꜓mi\\
  \begin{english}
    libero dalla sofferenza
  \end{english}

A꜕vero ho꜓mi\\
  \begin{english}
Libero da ostilità,
  \end{english}

A꜕byāpajjho ho꜓mi\\
  \begin{english} 
    Libero da malevolenza,
  \end{english}

A꜕nīgho ho꜓mi\\
  \begin{english}
Libero da ansia,
  \end{english}

Sukhī꜓ attānaṁ pa꜕riha꜓rāmi

  \begin{english}
E che possa mantenere il benessere in me stesso.
  \end{english}

\clearpage

Sa꜕bbe sa꜕ttā sukhitā ho꜓ntu\\
  \begin{english}
Che tutti gli esseri siano felici
  \end{english}

Sa꜕bbe sa꜕ttā averā ho꜓ntu\\
  \begin{english}
    Liberi da ostilità,
  \end{english}

Sa꜕bbe sa꜕ttā abyāpajjhā ho꜓ntu\\
  \begin{english}
Liberi da malevolenza,
  \end{english} 

Sa꜕bbe sa꜕ttā anīghā ho꜓ntu\\
  \begin{english}
Liberi da ansia,
  \end{english}

Sa꜕bbe sa꜕ttā sukhī꜓ a꜕ttānaṁ pa꜕riha꜓rantu
  \begin{english} 
E che possano mantenere il benessere in loro stessi.
  \end{english}

Sa꜕bbe sa꜕ttā sabbadukkhā pamucca꜓ntu
  \begin{english} 
che tutti gli esseri siano liberi dalla sofferenza.
  \end{english}

Sa꜕bbe sa꜕ttā laddha-sa꜓mpa꜕tti꜓to mā vigaccha꜓ntu

  \begin{english}
    e che non vengano separati dalla fortuna ottenuta.
  \end{english}

%\suttaref{AN 5.57}%
Sa꜕bbe sa꜕ttā kammassa꜕kā kamma꜓dāyādā kamma꜓yonī\\
  \begin{english}
tutti gli esseri sono padroni del loro kamma,\\
 eredi del loro kamma e nati dal loro kamma,
  \end{english}

kamma꜓bandhū kammapa꜕ṭisa꜓ra꜕ṇā\\
  \begin{english}
    legati al loro kamma e sostenuti del loro kamma,
  \end{english} 
    Yaṁ kammaṁ ka꜕rissa꜓nti\\

  \begin{english}
Qualsiasi kamma compiano,
  \end{english}

Kalyāṇaṁ vā pāpa꜕kaṁ vā\\

  \begin{english}
Buono o cattivo,
  \end{english}

Tassa꜕ dāyādā bha꜕vissa꜓nti

  \begin{english}
ne saranno eredi.
  \end{english}

\chapter[Dimore divine]{Diffusione delle dimore divine}

\firstline{Mettā-sahagatena}

\delegateSetUseNext

\begin{leader}
  [Ha꜓nda mayaṁ caturappamaññā-obhāsanaṁ karomase]
\end{leader}

\begin{english}
  [Facciamo ora risplendere i quattro immensurabili]
\end{english}

%\suttaref{DN 13}%
[Mettā-sa꜕ha꜕ga꜕tena] cetasā ekaṁ disaṁ pha꜕ri꜕tv꜕ā viha꜕ra꜕ti\\
  \begin{english} 
    [Dimorerò] pervadendo una direzione con il cuore imbevuto di gentilezza amorevole,
\end{english}

Ta꜕thā dutiyaṁ ta꜕thā tatiyaṁ ta꜕thā ca꜕tutthaṁ\\

\begin{english}
Così la seconda, così la terza, così la quarta,
\end{english}

Iti uddhamadho tiriyaṁ sabba꜕dhi꜕ sabbatta꜕tāya\\

\begin{english} 
  Così sopra e sotto,\\
   intorno e in ogni luogo,\\
   a tutti come a me stesso.
\end{english}

Sabbāvantaṁ lokaṁ mettā-sa꜕ha꜕ga꜕tena cetasā\\
 
\begin{english}
Dimorerò pervadendo l'intero mondo\\
 con un cuore imbevuto di gentilezza amorevole, 
  \end{english}


Vipulena mahagga꜕tena appa꜕māṇena a꜕verena a꜕byāpajjhena\\
pha꜕ri꜕tv꜕ā viha꜕ra꜕ti

  \begin{english}
abbondante, sublime, immensurabile,\\
senza ostilità e senza malevolenza.
  \end{english}

Karuṇā-sa꜕ha꜕ga꜕tena cetasā ekaṁ disaṁ pha꜕ri꜕tv꜕ā viha꜕ra꜕ti\\

\begin{english}
[Così dimorerò] pervadendo una direzione con il cuore imbevuto di compassione,
  \end{english}

Ta꜕thā dutiyaṁ ta꜕thā tatiyaṁ ta꜕thā ca꜕tutthaṁ\\

  \begin{english}
Così la seconda, così la terza, così la quarta,
  \end{english}

Iti uddhamadho tiriyaṁ sabba꜕dhi꜕ sabbatta꜕tāya\\

     \begin{english}
Così sopra e sotto,\\
   intorno e in ogni luogo,\\
   a tutti come a me stesso.
  \end{english}

Sabbāvantaṁ lokaṁ ka꜕ru꜕ṇā-sa꜕ha꜕ga꜕tena cetasā\\
  \begin{english} 
Dimorerò pervadendo l'intero mondo\\
 con un cuore imbevuto di compassione,
  \end{english}

Vipulena mahagga꜕tena appa꜕māṇena a꜕verena a꜕byāpajjhena\\
pha꜕ri꜕tv꜕ā viha꜕ra꜕ti
 
\begin{english} 
abbondante, sublime, immensurabile,\\ 
senza ostilità e senza malevolenza.
  \end{english} 

Muditā-sa꜕ha꜕ga꜕tena cetasā ekaṁ disaṁ pha꜕ri꜕tv꜕ā viha꜕ra꜕ti\\

\begin{english} 
Dimorerò pervadendo una direzione con il cuore imbevuto di gioia empatica,
  \end{english}

Ta꜕thā dutiyaṁ ta꜕thā tatiyaṁ ta꜕thā ca꜕tutthaṁ\\
\begin{english} 
Così la seconda, così la terza, così la quarta,
  \end{english} 

Iti uddhamadho tiriyaṁ sabba꜕dhi꜕ sabbatta꜕tāya\\

\begin{english} 
Così sopra e sotto,\\
   intorno e in ogni luogo,\\
   a tutti come a me stesso.
  \end{english}

Sabbāvantaṁ lokaṁ mu꜕di꜕tā-sa꜕ha꜕ga꜕tena cetasā\\

\begin{english} 
Dimorerò pervadendo l'intero mondo\\
 con un cuore imbevuto di gioia empatica,
  \end{english}


Vipulena mahagga꜕tena appa꜕māṇena a꜕verena a꜕byāpajjhena\\
pha꜕ri꜕tv꜕ā viha꜕ra꜕ti

  \begin{english} 
abbondante, sublime, immeasurabile,\\
senza ostilità e senza malevolenza.
  \end{english}


Upekkhā-saha꜕ga꜕te꜕na cetasā ekaṁ disaṁ pha꜕ri꜕tv꜕ā viha꜕ra꜕ti\\

\begin{english}
Dimorerò pervadendo una direzione con il cuore imbevuto di equanimità,
  \end{english}

Ta꜕thā dutiyaṁ ta꜕thā tatiyaṁ ta꜕thā ca꜕tutthaṁ\\

\begin{english}
Così la seconda, così la terza, così la quarta,
  \end{english}

  Iti uddhamadho tiriyaṁ sabba꜕dhi꜕ sabbatta꜕tāya\\

\begin{english} 
  Così sopra e sotto,\\
   intorno e in ogni luogo,\\
   a tutti come a me stesso.
  \end{english}

  Sabbāvantaṁ lokaṁ u꜕pe꜕kkhā-sa꜕ha꜕ga꜕tena cetasā\\

  \begin{english} 
Dimorerò pervadendo l'intero mondo\\  
con un cuore imbevuto di equanimità,
  \end{english}

  Vipulena mahagga꜕tena appa꜕māṇena a꜕verena a꜕byāpajjhena\\
  \vin pha꜕ri꜕tv꜕ā viha꜕ra꜕tī'ti

  \begin{english}
abbondante, sublime, immeasurabile,\\ 
senza ostilità e senza malevolenza.
  \end{english}


\chapter{Le più alte benedizioni}

\firstline{Thus have I heard that the Blessed One}

\begin{leader}
  [Cantiamo ora i versi delle più alte benedizioni]
\end{leader}

%\suttaref{Sn 2.4}%
[Così ho udito] il Beato una volta soggiornava a Sāvatthī,\\
 nella foresta di Jeta,\\ 
 nel parco di Anāthapiṇḍika.

Poi nel buio della notte,\\ 
 una divinità radiante illuminò tutta la foresta di Jeta.\\
Si inchinò al cospetto del Beato,\\
 quindi stando da un lato gli disse:

molti déi e uomini riflettono sulle massime benedizioni,\\
 desiderando felicità e sicurezza,\\
 quali sono dunque le più alte benedizioni?

Non associarsi agli stolti,\\
associarsi ai saggi,\\
venerare coloro che sono degni di lode,\\
queste sono le più alte benedizioni.

dimorare in luoghi adatti,\\
avere creato meriti in passato,\\
stabilire rettamente le proprie intenzioni,\\
queste sono le più alte benedizioni.

\clearpage

aver ampia conoscenza e molte abilità,\\
disciplina ben studiata e osservata,\\
parole ben dette, rette e armoniose,\\
queste sono le più alte benedizioni.

Assistere madre e padre,\\
Prendersi cura di figli e consorte,\\
Svolgere un lavoro non stressante,\\
Queste sono le più alte benedizioni.

Essere generosi e vivere rettamente,\\
Aiutare i propri parenti,\\
Compiere azioni irreprensibili,\\
Queste sono le più alte benedizioni.

Astenersi dal praticare il male,\\
Evitare sostanze che offuscano la mente,\\
Essere vigili in ogni circostanza,\\
Queste sono le più alte benedizioni.

Essere rispettosi e umili,\\
Essere contenti e grati,\\
Ascoltare il Dhamma nei momenti opportuni,\\
Queste sono le più alte benedizioni.

Essere pazienti e facili da ammonire,\\
Incontrare venerabili rinuncianti,\\
Parlare del Dhamma nei momenti opportuni,\\
Queste sono le più alte benedizioni.

\clearpage

Bruciare le contaminazioni e condurre la Vita Santa,\\
Vedere la Quattro Nobili Verità,\\
Realizzare il Nibbana,\\
Queste sono le più alte benedizioni.

Non essere scossi dalle vicissitudini del mondo,\\
Liberi da sofferenza, immacolati e sicuri,\\
Queste sono le più alte benedizioni.

Coloro che hanno compiuto questi atti vivono ovunque non sconfitti,\\
Ovunque son felici e sicuri,\\
Queste sono le più alte benedizioni.

\chapter[Nibbāna]{Riflessione sul Nibbāna}

\firstline{Atthi bhikkhave ajātaṁ abhūtaṁ akataṁ}

\begin{leader}
  [Ha꜓nda mayaṁ nibbāna-sutta-pāṭhaṁ bha꜕ṇāmase]
\end{leader}

\begin{english}
  [Cantiamo ora le parole del Buddha sul Nibbāna]
\end{english}


%\suttaref{Ud 8.3}%
Atthi bhi꜓kkha꜕ve a꜕jātaṁ a꜓bhūtaṁ a꜕kataṁ a꜕sa꜓ṅkh꜕ataṁ

\begin{english}
[C'è un Non-nato] Non-originato, Non-creato e Non-formato.
\end{english}

N꜕o cetaṁ bhi꜓kkha꜕ve a꜕bhavissa a꜕jātaṁ a꜓bhūtaṁ a꜕kataṁ a꜕sa꜓nkh꜕ataṁ

\begin{english}
Se non ci fosse questo Non-nato, Non-originato, Non-creato e Non-formato,
\end{english}

Na꜕ yidaṁ jātassa꜕ bhūtassa ka꜕tassa sa꜓ṅkh꜕atassa nissaraṇaṁ paññāye꜓tha

\begin{english}
la liberazione dal Non-nato, Non-originato, Non-creato e Non-formato non sarebbe possibile.
\end{english}

Ya꜕smā ca kho bhi꜓kkh꜕ave atthi a꜕jātaṁ a꜓bhūtaṁ a꜕kataṁ a꜕sa꜓ṅkha꜕taṁ

\begin{english}
Ma dato che c'è un Non-nato, Non-originato, Non-creato e Non-formato,
\end{english}

Ta꜕smā jātass꜕a bhūtassa ka꜕tassa sa꜓ṅkha꜕tassa nissaraṇaṁ paññāyati

\begin{english}
allora è possibile la liberazione dal mondo del nato, originato, creato e formato.
\end{english}

\clearpage

\chapter[Come i fiumi]{Come i fiumi}

\firstline{Yathā vāri-vahā pūrā paripūrenti sāgaraṁ}

%\suttaref{Khp 7.8}%
Yathā vāri-vahā꜓ pūrā pa꜕ripūrenti sāgaraṁ


\begin{english}
Cosi  come i fiumi pieni d'acqua riempiono completamente il mare
\end{english}

Evam-eva i꜓to dinnaṁ pe꜕tānaṁ u꜕pakappa꜕ti
 
\begin{english}
Così anche il merito donato ai defunti si accumula
\end{english}

%\suttaref{Dhp A I.198}%
Icchitaṁ pa꜕tthitaṁ tu꜓mhaṁ


\begin{english}
Che tutte le vostre speranze e tutti i vostri desideri
\end{english}

Khippam-eva sami꜓jjhatu

    
\begin{english}
Possano realizzarsi in breve tempo 
 \end{english}

Sa꜕bbe pūrentu sa꜓ṅkappā
 

\begin{english}
   Che tutti i vostri progetti si realizzino
\end{english}

Cando paṇṇaraso꜓ yathā

\begin{english}
Come la luna nel suo quindicesimo giorno
\end{english}

Maṇi jot꜕iraso꜓ yathā
 
\begin{english}
 O come una gemma splendente
\end{english}

%\suttaref{Trad.}%
Sabb'ītiyo vivajja꜓ntu


\begin{english}
che tutte le sventure siano evitate
\end{english}

S꜕abba-rogo vinassa꜕tu

\begin{english}
che tutte le malattie siano eliminate
\end{english}

Mā te bha꜕vatv-antarā꜓yo


\begin{english}
 possiate non incontrare mai pericoli
\end{english}

Sukhī꜓ dīgh'āyu꜕ko bha꜕va

\begin{english}
Possiate essere felici e vivere a lungo
\end{english}

%\suttaref{Dhp 109}%
A꜕bhivādana-sī꜓lissa꜕\\
Niccaṁ vu꜕ḍḍhāpa꜕cāyi꜕no\\
C꜕attāro dhammā vaḍḍha꜓nti\\
Āyu꜓ vaṇṇo su꜕khaṁ balaṁ
 
\begin{english}
  Con rispetto e buona condotta,\\
 onorando gli anziani,\\
  Quattro sono le qualità che aumentano:\\
  Vita, bellezza, felicità e forza.
\end{english}

%\suttaref{Trad.}%
Bhavatu sa꜕bba꜕-maṅg꜓alaṁ


\begin{english}
  Che ogni benedizione vi raggiunga
\end{english}

Rakkha꜓ntu sa꜕bba꜕-deva꜓tā

\begin{english}
  e tutti i Deva vi proteggano.
\end{english}

Sa꜕bba-bu꜓ddhānu꜓bhāvena

\begin{english}
  Attraverso il potere di tutti i Buddha
\end{english}

Sa꜕dā so꜕tthī꜓ bhavantu꜕ te

\begin{english}
Che voi siate sempre al sicuro.
\end{english}

\ifaivedition
\clearpage
\fi

Bhavatu sa꜕bba꜕-maṅg꜓alaṁ

\begin{english}
Che ogni benedizione vi raggiunga
\end{english}

Rakkha꜓ntu sa꜕bba꜕-deva꜓tā

\begin{english}
E tutti i Deva ti proteggano.
\end{english}

Sa꜕bba-dha꜓mmānu꜓bhāvena

\begin{english}
  Attraverso il potere di tutti i Dhamma
\end{english}

Sa꜕dā so꜕tthī꜓ bhavantu꜕ te

\begin{english}
Che voi siate sempre al sicuro.
\end{english}

Bhavatu sa꜕bba꜕-maṅg꜓alaṁ

\begin{english}
Che ogni benedizione vi raggiunga
\end{english}

Rakkha꜓ntu sa꜕bba꜕-deva꜓tā

\begin{english}
E tutti i Deva vi proteggano.
\end{english}

Sa꜕bba-sa꜓ṅghānu꜓bhāvena

\begin{english}
  Attraverso il potere di tutti i Sangha
\end{english}

Sa꜕dā so꜕tthī꜓ bhavantu꜕ te

\begin{english}
Che voi siate sempre al sicuro.
\end{english}

\chapter[I quattro requisiti]{Riflessione sui quattro requisiti}

\firstline{Paṭisaṅkhā yoniso}

\begin{leader}
  [Ha꜓nda mayaṁ taṅkhaṇika-paccave꜕kkhaṇa-pāṭhaṁ bhaṇāmase]
\end{leader}

\begin{english}
  [Cantiamo ora la riflessione sui quattro requisiti]
\end{english}

%\suttaref{MN 2}%
[Paṭisaṅkhā] yoniso cīva꜕raṁ pa꜕ṭise꜓vāmi, yāvadeva sī꜓tassa꜕\\
pa꜕ṭighātāya, uṇhassa pa꜕ṭighātāya, ḍaṁsa-maka꜕sa꜕-vātāta꜕pa꜕-siriṁsapa-\\
-samphassānaṁ pa꜕ṭighātāya, yāvadeva hiri꜓kopina-pa꜕ṭicchāda꜕natthaṁ

\begin{english}
[Riflettendo saggiamente] uso le vesti, solo per difendermi dal freddo, difendermi dal caldo, difendermi dal contatto di mosche, zanzare, vento, sole e crature striscianti, e solo per coprire il corpo.\end{english}

[Paṭisaṅkhā] yoniso piṇḍa꜕pātaṁ pa꜕ṭise꜓vāmi, neva da꜕vāya, na ma꜕dāya, na maṇḍa꜕nāya, na꜕ vi꜓bhūsa꜕nāya, yāvadeva i꜓massa꜕ kāyassa꜕ ṭhi꜕tiyā, yāpa꜕nāya, vihiṁsū꜕para꜓ti꜕yā, brahmaca꜕ri꜓yānugga꜕hāya, iti purāṇañca꜕ veda꜓naṁ pa꜕ṭiha꜓ṅkhāmi, navañca꜕ veda꜓naṁ na uppādessāmi, yātrā ca꜕ me bhavissati a꜕navajjatā ca꜕ phāsuvihāro cā'ti

\begin{english}
Rflettendo saggiamente uso il cibo offerto, non per piacere e divertimentimento, non per desiderio di euforia, non per ingrassare, non per ostentazione e vanità.
Lo uso solo per mantenere il corpo, per mantenerlo sano, per porre fine al disagio, per sostenere la Vita Santa,
alleviando la fame senza stramangiare continuerò a vivere senza colpa e mio agio.

\end{english}

[Paṭisaṅkhā] yoniso senāsa꜕naṁ pa꜕ṭise꜓vāmi, yāvadeva sī꜓tassa꜕\\
pa꜕ṭighātāya, uṇhassa pa꜕ṭighātāya, ḍaṁsa-maka꜕sa꜕-vātāta꜕pa꜕-siriṁsapa-\\
-samphassānaṁ pa꜕ṭighātāya, yāvadeva utupa꜕rissaya vi꜕nodanaṁ pa꜕ṭisa꜓llānārāmatthaṁ

\begin{english}
Riflettendo saggiamente uso la dimora, solo per difendermi dal freddo, difendermi dal caldo, difendermi dal contatto di mosche, zanzare, vento, sole e crature striscianti, solo per proteggermi dalle intemperie e per vivere in solitudine.
\end{english}

[Paṭisaṅkhā] yoniso gi꜕lāna-pacca꜕ya꜕-bhesajja-pa꜕rikkhāraṁ pa꜕ṭise꜓vāmi, yāvadeva uppa꜓nnānaṁ veyyābādhi꜕kānaṁ veda꜕nānaṁ pa꜕ṭighātāya, a꜕byāpajjha-pa꜕ramatāyā'ti

\begin{english}
Riflettendo saggiamente uso i rimedi curativi, solo per difendermi da sensazioni spiacevoli già sorte, per la massima libertà dalle malattie.
\end{english}

\chapter[Cinque temi]{Cinque temi per la frequente riflessione}

\firstline{Jarā-dhammomhi jaraṁ anatīto}

\begin{leader}
  %\suttaref{Cf. AN 5.57}%
  [Ha꜓nda mayaṁ abhiṇha-paccave꜕kkhaṇa-pāṭhaṁ bhaṇāmase]
\end{leader}

\begin{english}
  [Cantiamo ora i cinque temi per la frequente riflessione]
\end{english}

\sidepar{Uomini}%
[Jarā-dhammomhi꜕] jaraṁ a꜕na꜕tīto

\sidepar{Donne}%
[Jarā-dhammāmhi꜕] jaraṁ a꜕na꜕tītā

\begin{english}
Sono soggetto alla vecchiaia, non sono esente dalla vecchiaia.
\end{english}

\sidepar{m.}%
Byādhi꜓-dhammomhi꜕ byādhiṁ a꜕na꜕tīto

\sidepar{w.}%
Byādhi꜓-dhammāmhi꜕ byādhiṁ a꜕na꜕tītā

\begin{english}
Sono soggetto alla malattia, non sono esente dalla malattia.
\end{english}

\sidepar{m.}%
Ma꜕raṇa-dhammomhi꜕ ma꜕raṇaṁ a꜕na꜕tīto

\sidepar{w.}%
Ma꜕raṇa-dhammāmhi꜕ ma꜕raṇaṁ a꜕na꜕tītā

\begin{english}
Sono soggetto alla morte, non sono esente dalla morte.
\end{english}

Sa꜕bbehi me pi꜕yehi ma꜕nāpehi꜕ nānābhāvo vi꜕nābhāvo

\begin{english}
Tutto ciò che è mio, piacevole e amato, è soggetto al cambiamento, si separerà da me.
\end{english}

\sidepar{m.}%
Kammassa꜕komhi kamma꜓dāyādo kamma꜕yoni kamma꜕bandhu kammapa꜕ṭisa꜓ra꜕ṇo\\
Yaṁ kammaṁ ka꜕rissāmi, kalyāṇaṁ vā pāpa꜕kaṁ vā, tassa꜕ dāyādo bha꜕vissāmi

\clearpage

\sidepar{w.}%
Kammassa꜕kāmhi kamma꜓dāyādā kamma꜕yoni kamma꜕bandhu kammapa꜕ṭisa꜓ra꜕ṇā\\
Yaṁ kammaṁ ka꜕rissāmi, kalyāṇaṁ vā pāpa꜕kaṁ vā, tassa꜕ dāyādā bha꜕vissāmi

\begin{english}
Sono padrone del mio kamma, erede del mio kamma,\\
 nato dal mio kamma, legato al mio kamma, \\
 qualunque kamma compio, buono o cattivo ne sarò l'erede.
\end{english}

Evaṁ amhehi꜕ a꜕bhiṇhaṁ pacca꜕vekkhi꜓tabbaṁ


\begin{english}
  Così si dovrebbe riflettere frequentemente su questi cinque temi.
\end{english}

\chapter[Ten Subjects]{Ten Subjects for Frequent Recollection by One Who Has Gone Forth}

\firstline{Dasa ime bhikkhave}

\enlargethispage{\baselineskip}

\begin{leader}
  [Ha꜓nda mayaṁ pabbajita\hyp{}abhiṇha\hyp{}paccave꜕kkhaṇa\hyp{}pāṭhaṁ bhaṇāmase]
\end{leader}

%\suttaref{AN 10.48}%
[Dasa i꜕me bhikkhave] dhammā pabba꜕jitena a꜕bhiṇhaṁ pacca꜕vekkhi꜓tabbā, ka꜕ta꜕me dasa

\begin{english}
  Bhikkhus, there are te꜕n dhammas which should be re꜕flected upon again and a꜕gain by one who ha꜕s go꜓ne forth. \prul{What} a꜕re these ten?
\end{english}

Vevaṇṇi꜕yamhi ajjhūpa꜕ga꜕to'ti pabba꜕jitena a꜕bhiṇhaṁ pacca꜕vekkhi꜓tabbaṁ

\begin{english}
  `I am no꜕ longer li꜓ving according to꜕ worldly aims and va꜕lues.'\\
  This should be re꜕flected upon again and a꜕gain\\
  by one who ha꜕s go꜓ne forth.
\end{english}

Parapaṭi꜕baddhā me jīvi꜓kā'ti pabba꜕jitena a꜕bhiṇhaṁ pacca꜕vekkhi꜓tabbaṁ

\begin{english}
  `My very꜕ life is susta꜓ined through the gifts of o꜕thers.'\\
  This should be re꜕flected upon again and a꜕gain\\
  by one who ha꜕s go꜓ne forth.
\end{english}

Añño me ākappo ka꜕ra꜕ṇīyo'ti pabba꜕jitena a꜕bhiṇhaṁ pacca꜕vekkhi꜓tabbaṁ

\begin{english}
  `I shou꜕ld strive to aba꜓ndon my former ha꜕bits.'\\
  This should be re꜕flected upon again and a꜕gain\\
  by one who ha꜕s go꜓ne forth.
\end{english}

\clearpage

Kacci nu꜕ kho me attā sīla꜕to na u꜕pavadatī'ti pabba꜕jitena a꜕bhiṇhaṁ pacca꜕vekkhi꜓tabbaṁ

\begin{english}
  `Does re꜕gret over my co꜓nduct arise in m꜕y mind?'\\
  This should be re꜕flected upon again and a꜕gain\\
  by one who ha꜕s go꜓ne forth.
\end{english}

Kacci nu꜕ kho maṁ a꜕nuvicca viññū sabrahma꜓cārī sīla꜕to na u꜕pavadantī'ti pabba꜕jitena a꜕bhiṇhaṁ pacca꜕vekkhi꜓tabbaṁ

\begin{english}
  `Could m꜕y spiritual compa꜓nions find fault with my co꜕nduct?'\\
  This should be re꜕flected upon again and a꜕gain\\
  by one who ha꜕s go꜓ne forth.
\end{english}

Sa꜕bbehi me pi꜕yehi ma꜕nāpehi꜕ nānābhāvo vi꜕nābhāvo'ti pabba꜕jitena abhiṇhaṁ pacca꜕vekkhi꜓tabbaṁ

\begin{english}
  `All that i꜕s mine, be꜕loved and ple꜓asing, will become o꜕therwise, will become se꜓parated from me.'\\
  This should be re꜕flected upon again and a꜕gain\\
  by one who ha꜕s go꜓ne forth.
\end{english}

Kammassa꜕komhi kamma꜓dāyādo kamma꜕yoni kamma꜕bandhu kammapa꜕ṭisa꜓raṇo, yaṁ kammaṁ ka꜕rissāmi, kalyāṇaṁ vā pāpa꜕kaṁ vā, tassa꜕ dāyādo bha꜕vissāmī'ti pabba꜕jitena a꜕bhiṇhaṁ pacca꜕vekkhi꜓tabbaṁ

\begin{english}
  `I am the꜕ owner of my ka꜕mma, heir to my ka꜕mma, born of my ka꜕mma,\\
  re꜕lated to my ka꜕mma, abide suppo꜓rted by my ka꜕mma; whatever kamma I sha꜕ll do, for good or fo꜕r ill, of \prul{that} I will be꜕ the꜓ heir.'\\
  This should be re꜕flected upon again and a꜕gain\\
  by one who ha꜕s go꜓ne forth.
\end{english}

\clearpage

`Kathambhūtassa꜕ me rattindi꜕vā vīti꜕pa꜓tantī'ti pabba꜕jitena a꜕bhiṇhaṁ pacca꜕vekkhi꜓tabbaṁ

\begin{english}
  `The꜕ days and nights are re꜕lentlessly pa꜓ssing; ho꜕w well am I spe꜓ndi꜓ng\\ m꜕y time?'\\
  This should be re꜕flected upon again and a꜕gain\\
  by one who ha꜕s go꜓ne forth.
\end{english}

Kacci nu꜕ kho'haṁ suññā꜓gāre abhira꜕māmī'ti pabba꜕jitena a꜕bhiṇhaṁ pacca꜕vekkhi꜓tabbaṁ

\begin{english}
  `Do I delight in so꜓litude or not?'\\
  This should be re꜕flected upon again and a꜕gain\\
  by one who ha꜕s go꜓ne forth.
\end{english}

Atthi nu꜕ kho me uttari-ma꜕nussa-dhammā alamariya꜕-ñāṇa-dassana-viseso adhiga꜕to, so'haṁ pacchi꜓me kāle sa꜕brahmacārīhi꜕ puṭṭho na maṅku bha꜕vissāmī'ti pabba꜕jitena a꜕bhiṇhaṁ pacca꜕vekkhi꜓tabbaṁ

\begin{english}
  `Has m꜕y practice borne fruit with freedom or i꜓nsight so that at the e꜕nd of my life I need not feel a꜕shamed when questioned b꜕y my spi꜓ri꜓tual compa꜕nions?'\\
  This should be re꜕flected upon again and a꜕gain\\
  by one who ha꜕s go꜓ne forth.
\end{english}

Ime kho bhikkha꜓ve da꜕sa꜕ dhammā pabba꜕jitena a꜕bhiṇhaṁ pacca꜕vekkhitabbā'ti

\begin{english}
  Bhikkhus, these are the te꜕n dhammas to be re꜕flected upon again and a꜕gain by one who ha꜕s go꜓ne forth.
\end{english}

\chapter[Thirty-Two Parts]{Reflection on the Thirty-Two Parts}

\firstline{Ayaṁ kho me kāyo}

\begin{leader}
  [Ha꜓nda mayaṁ dvattiṁsākāra-pāṭhaṁ bhaṇāmase]
\end{leader}

%\suttaref{DN 22}%
[Ayaṁ kho] me kāyo uddhaṁ pāda꜕ta꜕lā adho kesamatthakā\\
ta꜕ca꜕pa꜕ri꜕yanto pūro nānappa꜕kārassa꜕ a꜕su꜕ci꜕no

\begin{english}
  This, which is my body, from the soles of the feet up, and down from the crown of the head, is a sealed bag of skin filled with unattractive things.
\end{english}

Atthi imasmiṁ kāye

\begin{english}
  In this body there are:
\end{english}

{\centering
\setArrayStrech{1}

\begin{tabular}{ r l }
kesā            & \tr{hair of the head} \\
lomā            & \tr{hair of the body} \\
nakhā           & \tr{nails} \\
dantā           & \tr{teeth} \\
taco            & \tr{skin} \\
maṁsaṁ          & \tr{flesh} \\
nahārū          & \tr{sinews} \\
aṭṭhī           & \tr{bones} \\
aṭṭhimiñjaṁ     & \tr{bone marrow} \\
vakkaṁ          & \tr{kidneys} \\
hadayaṁ         & \tr{heart} \\
yakanaṁ         & \tr{liver} \\
kilomakaṁ       & \tr{membranes} \\
pihakaṁ         & \tr{spleen} \\
papphāsaṁ       & \tr{lungs} \\
\end{tabular}

\clearpage

\begin{tabular}{ r l }
antaṁ           & \tr{bowels} \\
antaguṇaṁ       & \tr{entrails} \\
udariyaṁ        & \tr{undigested food} \\
karīsaṁ         & \tr{excrement} \\
pittaṁ          & \tr{bile} \\
semhaṁ          & \tr{phlegm} \\
pubbo           & \tr{pus} \\
lohitaṁ         & \tr{blood} \\
sedo            & \tr{sweat} \\
medo            & \tr{fat} \\
assu            & \tr{tears} \\
vasā            & \tr{grease} \\
kheḷo           & \tr{spittle} \\
siṅghāṇikā      & \tr{mucus} \\
lasikā          & \tr{oil of the joints} \\
muttaṁ          & \tr{urine} \\
matthaluṅgan'ti & \tr{brain} \\
\end{tabular}

\restoreArrayStretch
}

Evam-ayaṁ me kāyo uddhaṁ pāda꜕ta꜕lā adho kesamatthakā\\
ta꜕ca꜕pa꜕ri꜕yanto pūro nānappa꜕kārassa꜕ a꜕su꜕ci꜕no

\begin{english}
  This, then, which is my body, from the soles of the feet up, and down from the crown of the head, is a sealed bag of skin filled with unattractive things.
\end{english}

% End of reflections-and-recollections-p1.tex
